\subsection{System Requirements and Specification}

This section summarizes the requirements and constraints that govern the design
of the embedded DSP Eurorack synthesizer module. The full requirements and
technical specification are documented separately and form the contractual basis
for implementation and acceptance.

\subsubsection{Functional Requirements}

The system shall acquire audio signals via an audio codec, process them
in real time using DSP algorithms executing on the MCU, and output the processed
signal through the same codec. The following functional capabilities are required:

\begin{itemize}
    \item Stereo audio input and output via 3.5\,mm mono jack sockets
    \item Real-time DSP processing with deterministic latency
    \item Three control voltage (CV) inputs, one V/Oct pitch input, and four gate
    inputs for Eurorack-compatible modulation, also connectable via 3.5 mm jack sockets
    \item Four slide potentiometers and two push buttons for direct user control
    of DSP parameters and playback modes
    \item Four WS2812 RGB LEDs for visual feedback
    \item Concurrent handling of audio processing and control interface updates
    without mutual interference
\end{itemize}

\subsubsection{Electrical Constraints}

The hardware must be compatible with the Eurorack modular synthesizer standard,
which imposes specific electrical requirements on signal levels and mechanical
dimensions.

\begin{itemize}
    \item All inputs must tolerate the full Eurorack signal range of approximately
    $-10$\,V to $+10$\,V without damage to downstream circuitry
    \item Gate inputs are protected via transistor clamping circuits; CV inputs
    use rail-guarded op-amp stages to prevent overvoltage at the MCU ADC
    pins
    \item Audio inputs are attenuated, diode-protected, and AC-coupled to
    safeguard the codec's analog front end
    \item Power is derived from the Eurorack bus; a 5\,V buck converter supplies
    the WS2812 LEDs, while a 3.3\,V LDO provides a clean supply for the
    MCU and analog circuitry
    \item PCB layout must minimize noise coupling between digital and analog
    signal paths
    \item Mechanical dimensions conform to the Eurorack 3U standard as derived
    from \textbf{IEC~60297} / \textbf{DIN~41494}
\end{itemize}

\subsubsection{Firmware Constraints}

The firmware must satisfy real-time processing requirements while maintaining a
modular architecture that supports future extension to other MCU-based Eurorack
modules.

\begin{itemize}
    \item Audio acquisition and output are driven by DMA transfers to ensure
    uninterrupted data flow independent of task scheduling
    \item FreeRTOS provides preemptive task scheduling; high-priority audio
    processing tasks are strictly separated from lower-priority control and
    interface tasks
    \item DSP algorithms are implemented in C, with CMSIS-DSP library routines
    used where applicable to exploit the STM32H7's hardware acceleration
    \item The audio codec (TLV320AIC3204) is configured at startup via a custom
    I2C driver; its digital audio stream is transferred via I2S
    \item Calibration data may be persisted to the MCU's internal flash;
    no external mass storage is required
\end{itemize}